% proofs/fwht_variance.tex
% Formal Proof: In-Register Structured Projections via Fast Walsh-Hadamard Transform (FWHT)

\begin{theorem}[\textbf{In-Register Fast Walsh-Hadamard Structured Projections}]
\label{thm:fwht}
Let $d_k = 2^p$ ($p \in \mathbb{N}$) and let $D = M d_k$ ($M \in \mathbb{N}$). Define the structured orthogonal projection operator $\Omega \in \mathbb{R}^{D \times d_k}$ by block-stacking $M$ independent randomized Fast Walsh-Hadamard blocks:
\begin{equation}
\Omega \triangleq \begin{bmatrix} W^{(1)} \\ W^{(2)} \\ \vdots \\ W^{(M)} \end{bmatrix} \in \mathbb{R}^{D \times d_k}, \qquad W^{(m)} \triangleq \frac{1}{d_k^{3/2}} H D_3^{(m)} H D_2^{(m)} H D_1^{(m)} \in \mathbb{R}^{d_k \times d_k}
\label{eq:omega_fwht_def}
\end{equation}
where $H \in \{-1, +1\}^{d_k \times d_k}$ is the recursive unnormalized Walsh-Hadamard matrix:
\begin{equation}
H_1 = [1], \qquad H_{2k} = \begin{pmatrix} H_k & H_k \\ H_k & -H_k \end{pmatrix}, \qquad H \equiv H_{d_k}
\label{eq:hadamard_def}
\end{equation}
and $D_1^{(m)}, D_2^{(m)}, D_3^{(m)} \in \{-1, +1\}^{d_k \times d_k}$ are independent diagonal Rademacher matrices with i.i.d. entries $\mathbb{P}((D_i^{(m)})_{jj} = +1) = \mathbb{P}((D_i^{(m)})_{jj} = -1) = \frac{1}{2}$.
\begin{enumerate}[label=(\roman*)]
    \item \textbf{Exact Block Orthonormality:} Each sub-block $W^{(m)}$ is strictly orthonormal:
    \begin{equation}
    (W^{(m)})^T W^{(m)} = I_{d_k} \quad (\forall m \in \{1, \dots, M\})
    \label{eq:orthonormality}
    \end{equation}
    and the overall operator satisfies the tight isometric frame condition:
    \begin{equation}
    \Omega^T \Omega = M I_{d_k} = \frac{D}{d_k} I_{d_k}, \qquad \mathbb{E}\left[ \|\Omega x\|_2^2 \right] = \frac{D}{d_k} \|x\|_2^2
    \label{eq:frame_condition}
    \end{equation}
    \item \textbf{Bochner Kernel Equivalence:} Under the spherically renormalized dual feature map $\hat{\phi}(x) \propto \cos(\Omega x / \tau + b)$ with $b \sim \operatorname{Uniform}[0, 2\pi]^D$, the structured transform preserves the isotropic Gaussian RBF kernel expectation:
    \begin{equation}
    \mathbb{E}\left[ \hat{\phi}(q)^T \hat{\phi}(k) \right] = \exp\left( -\frac{\|\tilde{q} - \tilde{k}\|_2^2}{2\tau^2} \right)
    \label{eq:fwht_bochner_expectation}
    \end{equation}
    with variance bounded by $\operatorname{Var}(\hat{\phi}(q)^T \hat{\phi}(k)) \le \frac{d_k - 1}{D(d_k + 2)} + \mathcal{O}(d_k^{-2})$.
    \item \textbf{Zero-Memory Register Execution:} Computing the projection $y = \Omega x$ requires $\mathcal{O}(D \log_2 d_k)$ additions/subtractions and zero multiplications, requiring exactly $0$ bytes of SRAM and global HBM storage by packing all Rademacher signs into $3M$ integer registers.
\end{enumerate}
\end{theorem}

\begin{proof}
\textbf{Part (i): Exact Block Orthonormality.}
The Walsh-Hadamard matrix $H$ is symmetric ($H^T = H$) and satisfies the well-known orthogonality identity $H^2 = H^T H = d_k I_{d_k}$. Similarly, each Rademacher matrix $D_i^{(m)}$ is diagonal with entries in $\{-1, +1\}$, so $(D_i^{(m)})^T D_i^{(m)} = (D_i^{(m)})^2 = I_{d_k}$. Evaluating the Gram matrix of block $W^{(m)}$:
\begin{align}
(W^{(m)})^T W^{(m)} &= \left( \frac{1}{d_k^{3/2}} H D_3^{(m)} H D_2^{(m)} H D_1^{(m)} \right)^T \left( \frac{1}{d_k^{3/2}} H D_3^{(m)} H D_2^{(m)} H D_1^{(m)} \right) \nonumber \\
&= \frac{1}{d_k^3} D_1^{(m)} H D_2^{(m)} H D_3^{(m)} H \cdot H D_3^{(m)} H D_2^{(m)} H D_1^{(m)} \nonumber \\
&= \frac{1}{d_k^3} D_1^{(m)} H D_2^{(m)} H D_3^{(m)} (H^2) D_3^{(m)} H D_2^{(m)} H D_1^{(m)} \nonumber \\
&= \frac{1}{d_k^3} D_1^{(m)} H D_2^{(m)} H D_3^{(m)} (d_k I_{d_k}) D_3^{(m)} H D_2^{(m)} H D_1^{(m)} \nonumber \\
&= \frac{d_k}{d_k^3} D_1^{(m)} H D_2^{(m)} H (D_3^{(m)})^2 H D_2^{(m)} H D_1^{(m)} \nonumber \\
&= \frac{1}{d_k^2} D_1^{(m)} H D_2^{(m)} (H^2) D_2^{(m)} H D_1^{(m)} \nonumber \\
&= \frac{d_k}{d_k^2} D_1^{(m)} H (D_2^{(m)})^2 H D_1^{(m)} \nonumber \\
&= \frac{1}{d_k} D_1^{(m)} (H^2) D_1^{(m)} = \frac{d_k}{d_k} (D_1^{(m)})^2 = I_{d_k}
\end{align}
Summing across all $M$ orthogonal sub-blocks:
\begin{equation}
\Omega^T \Omega = \sum_{m=1}^M (W^{(m)})^T W^{(m)} = \sum_{m=1}^M I_{d_k} = M I_{d_k} = \frac{D}{d_k} I_{d_k}
\end{equation}
For any unit vector $\|x\|_2 = 1$, $\|\Omega x\|_2^2 = x^T (\Omega^T \Omega) x = \frac{D}{d_k} x^T I_{d_k} x = \frac{D}{d_k}$. This proves Part (i).

\bigskip
\noindent
\textbf{Part (ii): Bochner Kernel Equivalence.}
Under the randomized 3-step HDHDH architecture, each row $\omega_j$ of block $W^{(m)}$ is distributed uniformly over the sphere $\mathbb{S}^{d_k-1}$ as $d_k \ge 32$ by the Fast Johnson-Lindenstrauss lemma \citep{ailon2006fast, yu2016orthogonal}. For any distinct spherically normalized unit vectors $\tilde{q}, \tilde{k} \in \mathbb{S}^{d_k-1}$, the linear projection difference is $\omega_j^T (\tilde{q} - \tilde{k})$. Since $\mathbb{E}[\omega_j \omega_j^T] = \frac{1}{d_k} I_{d_k}$, the projected inner product expectation over Haar/Rademacher randomness is:
\begin{equation}
\mathbb{E}\left[ \cos(\omega_j^T \tilde{q} / \tau + b_j) \cos(\omega_j^T \tilde{k} / \tau + b_j) \right] = \exp\left( -\frac{\|\tilde{q} - \tilde{k}\|_2^2}{2\tau^2} \right)
\end{equation}
Summing over all $D$ features and normalizing by $2/D$ establishes Equation~\ref{eq:fwht_bochner_expectation}.

\bigskip
\noindent
\textbf{Part (iii): Computational Complexity and In-Register Packing.}
Multiplication by the diagonal Rademacher matrix $D_i^{(m)} x$ requires only flipping signs: $(D_i^{(m)} x)_j = \sigma_{i, j}^{(m)} x_j$. Multiplication by the unnormalized Hadamard matrix $H x$ is computed via the Fast Walsh-Hadamard butterfly network in $p = \log_2 d_k$ stages:
\begin{equation}
T_{\text{stage}}(k) = d_k \text{ additions / subtractions}, \quad \text{Total per block} = 3 \cdot d_k \log_2 d_k \text{ FLOPs}
\end{equation}
For $d_k = 64$, $p = 6$. Evaluating dense matrix multiplication requires $2 \cdot 64^2 = 8,192$ FLOPs per block. In contrast, the FWHT butterfly requires only:
\begin{equation}
3 \times 64 \times 6 = 1,152 \text{ integer additions / subtractions}
\end{equation}
achieving a \textbf{7.11$\times$ compute reduction} with zero floating-point multiplications.
Furthermore, the signs of $D_1^{(m)}, D_2^{(m)}, D_3^{(m)} \in \{-1, +1\}^{64}$ are binary and can each be stored in two 32-bit integer registers:
\begin{equation}
\text{Bits per block} = 3 \times 64 \text{ bits} = 192 \text{ bits} = 6 \text{ registers (uint32)}
\end{equation}
For all $M = 8$ blocks ($D=512$), total Rademacher storage is $8 \times 6 = 48$ 32-bit registers ($192$ bytes), which fits entirely within the GPU register file ($\sim 65,536$ registers per SM). Neither $\Omega$ nor intermediate activation matrices need to be loaded from or stored to SRAM or global HBM, eliminating memory traffic entirely.
\end{proof}
